\documentclass[10pt,a4paper]{amsart}
\usepackage[margin=19mm]{geometry}
\usepackage{amsmath,amssymb,mathtools}
\usepackage[scr=rsfs,cal=euler]{mathalfa}
\usepackage{xfrac}
\usepackage[unicode,hidelinks,bookmarks=false]{hyperref}
\newcommand{\ie}{i.e.\ }
\newcommand{\cf}{cf.\ }
\newcommand{\resp}{respectively\ }
\newcommand{\Subsection}{\subsection}
\providecommand{\nobreakdash}{\nobreak\hbox{-}}
\setlength{\emergencystretch}{2em}
\multlinegap0pt
\usepackage{bm}
\usepackage{bbm}
\usepackage{dsfont}
\usepackage{xspace}
\usepackage{cancel}
\usepackage{mathtools}
\usepackage{amsthm}
\makeatletter
\@ifundefined{theorem}{\newtheorem{theorem}{Theorem}}{}
\@ifundefined{lemma}{\newtheorem{lemma}[theorem]{Lemma}}{}
\@ifundefined{prop}{\newtheorem{prop}[theorem]{Proposition}}{}
\makeatother
\DeclareMathOperator{\Gal}{Gal}
\newcommand{\A}{\mathbb{A}}
\newcommand{\Frob}{\mathrm{Frob}}
\newcommand{\F}{\mathbb{F}}
\newcommand{\GO}{\mathscr{G\!O}}
\newcommand{\Ord}{\mathrm{Ord}} %ordinary primes
\newcommand{\Q}{\mathbb{Q}}
\newcommand{\Z}{\mathbb{Z}}
\newcommand{\cC}{\mathcal{C}}
\newcommand{\cK}{\mathcal{K}}
\newcommand{\cO}{\mathcal{O}}
\newcommand{\cP}{\mathcal{P}}
\newcommand{\cS}{\mathcal{S}}
\newcommand{\fP}{\mathfrak{P}}
\newcommand{\fp}{\mathfrak{p}}
\newcommand{\vA}{{\vec{A}}}
\newcommand{\va}{{\vec{a}}}
\renewcommand{\bar}{\overline}
\title[Generically Ordinary One-parameter Hyperelliptic Families ar]{Generically Ordinary One-parameter Hyperelliptic Families are Dense}
\author{Hui June Zhu}
\date{Source: arXiv:2608.22779v1 (CC BY 4.0)}
\begin{document}
\maketitle

\noindent\emph{Source and licence.} Generically Ordinary One-parameter Hyperelliptic Families are Dense. Version arXiv:2608.22779v1 (CC BY 4.0), distributed under the Creative Commons Attribution 4.0 International licence, as recorded in the arXiv licence metadata for this exact version and shown on the abstract page https://arxiv.org/abs/2608.22779v1. Full rights notice: licenses/arxiv-2608.22779-CC-BY-4.0.txt.\par\smallskip

\noindent\textit{Editorial scope.} Counted unit: the Theorem whose source label is \texttt{None} in the article (source0.tex lines 1404--1428, source label \texttt{None}), delivered together with the article's own complete original proof of it (source0.tex lines 1430--1552, one \texttt{proof} environment of 123 lines). No mathematics is written, repaired or re-derived: statement and proof are the article's own text line for line. Required context retained uncounted: T:main (lines 198--228); T:main3 (lines 258--275); T:density (lines 291--307); P:Delta (lines 702--709); L:degree1 (lines 1387--1390). Every definition, display and label of those lines is retained unchanged. Omitted from this package and not redistributed: 1--197, 229--257, 276--290, 308--701, 710--1386, 1391--1403, 1429--1429, 1553--1583. Nothing inside a retained excerpt is omitted or altered: every retained body is delivered exactly as the article's lines are written, so no omission receipt of any kind accompanies this package. Citations: the retained excerpts carry no citation command at all, so no bibliography entry of the article is transcribed and no reference is redistributed. Reference adaptation: none was needed and none was made; every reference inside the delivered unit resolves inside it, so no unresolved reference or citation is produced. Printed slips kept literal: no printed word, symbol, hypothesis or number of the counted statement or of its proof was altered. Layout and class: the article's own class is \texttt{amsart} with packages including amsmath, hyperref, amsfonts, amssymb, amsthm, mathrsfs, mathtools, comment, pdfpages, inputenc, graphicx, fontenc; the wrapper is the lane's pinned \texttt{amsart} editorial wrapper at 10pt on A4, which restates the theorem-like environments the retained text needs and adds the article's own macro definitions that the retained text needs (\texttt{\textbackslash A}, \texttt{\textbackslash F}, \texttt{\textbackslash Frob}, \texttt{\textbackslash GO}, \texttt{\textbackslash Gal}, \texttt{\textbackslash Ord}, \texttt{\textbackslash Q}, \texttt{\textbackslash Z}, \texttt{\textbackslash bar}, \texttt{\textbackslash cC}, \texttt{\textbackslash cK}, \texttt{\textbackslash cO}, \texttt{\textbackslash cP}, \texttt{\textbackslash cS}, \texttt{\textbackslash fP}, \texttt{\textbackslash fp}, \texttt{\textbackslash vA}, \texttt{\textbackslash va}), with extra wrapper packages (bm, bbm, dsfont, xspace, cancel, mathtools). The delivered numbering is the unit's own. Assets: no retained span contains a figure environment, a graphics-inclusion command, a PDF-page inclusion, a table or a TikZ picture; no image file is copied into this package, the package declares no asset dependency and no image file is redistributed with it. Licence: version arXiv:2608.22779v1 of this article is distributed under the Creative Commons Attribution 4.0 International licence, as recorded in the arXiv licence metadata for this exact version and shown on the abstract page https://arxiv.org/abs/2608.22779v1. A full rights notice is delivered at licenses/arxiv-2608.22779-CC-BY-4.0.txt. Characters: every retained excerpt is pure ASCII, so the delivered engine, XeLaTeX with its default Latin Modern text font, reports no missing character.
\begin{theorem}\label{T:main}
Let $d=2g+1$ with $g\ge 2$.
There exist Zariski-open dense subschemes
$U_k\subseteq \A_\Z^{d-1}$ for $k\in(\Z/d\Z)^\times$ such that
$U\coloneqq \bigcap_k U_k$ is Zariski-open and dense, and the following statements hold.
\begin{enumerate}
\item Let
$\va\in U_\Q(\bar\Q)\cap\bar\Z^{d-1}$ and $F=\Q(\va)$. Then
there exists $B(d,\va)>0$
such that,
for every nonzero prime
ideal $\fp\subset\cO_F$ whose residue characteristic $p$ satisfies
$p>B(d,\va)$,
the reduction of
$
\cC_{\va}\longrightarrow \cS_\va^\circ
$
modulo $\fp$
is generically ordinary.
\item
Let $k\in(\Z/d\Z)^\times$,
$\va\in U_{k,\Q}(\bar\Q)\cap {\bar\Z}^{d-1}$, and $F=\Q(\va)$.
Then there exists $B_k(d,\va)>0$
such that,
for every nonzero prime ideal $\fp\subset \cO_F$
whose residue characteristic $p$ satisfies $p\equiv k\pmod d$
and $p>B_k(d,\va)$, the reduction of
$\cC_{\va}\longrightarrow \cS_\va^\circ$
modulo $\fp$ is generically ordinary.
\end{enumerate}
\end{theorem}

\begin{theorem}\label{T:main3}
Let $d=2g+1$ with $g\ge 2$.
\begin{enumerate}
\item Let $\va=(a_1,\ldots,a_{d-1})\in\bar\Z^{d-1}$ be an arbitrary coefficient vector.
Let $L$ be a finite extension of $\Q(\va)$ and let
$\cC_{\va,\cO_L}\longrightarrow \cS_{\va,\cO_L}^\circ$
denote the base change to $\cO_L$.
Then for every nonzero prime ideal $\fP\subset \cO_L$ whose residue characteristic $p$ satisfies $p\equiv 1\pmod d$,
the reduction of $\cC_{\va,\cO_L}\longrightarrow\cS_{\va,\cO_L}^\circ$
modulo $\fP$ is generically ordinary.

\item For every prime $p\equiv 1\pmod d$
and every coefficient vector $\va\in\bar\F_p^{d-1}$,
the corresponding hyperelliptic family
is generically ordinary.
That is, $\GO_p=\A_{\F_p}^{d-1}$.
\end{enumerate}
\end{theorem}

\begin{theorem}\label{T:density}
Let $d=2g+1$ with $g\ge 2$. Let $\zeta_d$ denote a primitive
$d$-th root of unity in $\bar\Q$.
Let $\va=(a_1,\ldots,a_{d-1})\in{\bar\Z}^{d-1}$, put $F=\Q(\va)$, and
$L=F(\zeta_d)$.
Let $C$ be the hyperelliptic curve
given by the affine equation
\[
y^2=f_{\va,t}(x)
\]
over $F(t)$.
Then the set of ordinary primes of $C/F(t)$
has natural lower density at least $\frac{1}{[L:F]}$.
If $C_L\coloneqq C\times_{F(t)} L(t)$,
then the set of ordinary primes of $C_L/L(t)$ has
natural density~$1$.
\end{theorem}

\begin{prop}\label{P:Delta}
Let $\Delta_k^\Z(\vA)=\alpha^g\Delta_k(\vA)$.
Then $\Delta_k^\Z(\vA)\in\Z[\vA]$ is nonzero.
For any prime $p>d$,
the reductions of $\Delta_k\in\Z_{(p)}[\vA]$ and $\Delta_k^\Z\in\Z[\vA]$
modulo $p$ are nonzero.
Moreover, $\Delta_1=1$.
\end{prop}

\begin{lemma}\label{L:degree1}
For any number field $L$, the density of $\cP_{1,L}$ is $1$,
and that of $\cP_{2^+,L}$ is $0$.
\end{lemma}

\begin{theorem}[A refinement of Theorem~\ref{T:density}]
Let $d=2g+1$ with $g\ge 2$.
Let $\va=(a_1,\ldots,a_{d-1})\in{\bar\Z}^{d-1}$
and $F=\Q(\va)$.
Let $C/F(t)$ be the smooth projective hyperelliptic curve
given by the affine equation
\[
y^2=x^d+a_{d-1}x^{d-1}+\cdots+a_1x+t.
\]
Let $L=F(\zeta_d)$, and identify
 $H=\Gal(L/F)$ with a subgroup of $(\Z/d\Z)^\times$ via its action on $\zeta_d$, and
put
\[
\cK_\va\coloneqq \{k\in H:\Delta_k^\Z(\va)\ne 0\}.
\]
\begin{enumerate}
\item
The lower density of ordinary primes for $C/F(t)$ satisfies
\[
\underline\delta_F(\Ord(C/F(t)))\ge \frac{|\cK_\va|}{[L:F]}\ge \frac{1}{[L:F]}.
\]
\item
Let $C_L\coloneqq C\times_{F(t)}L(t)$. Then the density of ordinary primes for $C_L/L(t)$ is $1$.
\end{enumerate}
\end{theorem}

\begin{proof}
(1)
Let $\va\in\bar\Z^{d-1}$.
From Proposition~\ref{P:Delta} we have $\Delta_1=1$, and hence
\[
\Delta_1^\Z=\alpha^g\Delta_1=\alpha^g\ne 0.
\]
Therefore $1\in\cK_\va$, so in particular $|\cK_\va|\ge 1$.

For every $k\in\cK_\va$, we have $\va\in U_{k,\Q}(\bar\Q)\cap\bar\Z^{d-1}$.
By Theorem~\ref{T:main}(2), if we write $p$ for the rational prime lying under $\fp$, then
\begin{equation}\label{E:fp}
\{\fp\in\Pi_F:p\equiv k\pmod d\}\backslash\Ord(C/F(t))
\text{ is finite.}
\end{equation}
Write
\[
\cP_{1,F,k}\coloneqq \{\fp\in\cP_{1,F}:p\equiv k\pmod d\}.
\]
Then \eqref{E:fp} gives
$\cP_{1,F,k}\backslash\Ord(C/F(t))$ is finite. That is
\begin{equation}\label{E:finite}
\cP_{1,F,k}\text{ and }\Ord(C/F(t))\cap\cP_{1,F,k}
\text{ differ by a finite set}.
\end{equation}
Since $L=F(\zeta_d)$, the cyclotomic extension $L/F$ is unramified away from
the primes above $d$. Set
\[
\cP_F^{(d)}\coloneqq \{\fp\in\Pi_F: \fp\mid d\cO_F\}.
\]
This is a finite set containing every prime that ramifies in $L/F$.
For every prime $\fp\in \cP_{1,F}\backslash \cP_F^{(d)}$ we have
$N\fp=p$, and the arithmetic Frobenius $\Frob_\fp\in H$ acts as
\[
\Frob_\fp(\zeta_d)=\zeta_d^{N\fp}=\zeta_d^p.
\]
Under the identification $H\hookrightarrow (\Z/d\Z)^\times$,
it follows that, for every $k\in H$,
\[
{\Frob}_\fp=k \quad\Longleftrightarrow \quad \zeta_d^p=\zeta_d^k \quad
\Longleftrightarrow
\quad p\equiv k\pmod d.
\]
Consequently,
\[
\cP_{1,F,k}\backslash\cP_F^{(d)}
=(\cP_{1,F}\backslash\cP_F^{(d)})
\cap\{\fp\in\Pi_F\backslash\cP_F^{(d)}:\Frob_\fp=k\}.
\]
Then
\begin{equation}\label{E:delta1}
\cP_{1,F,k}\text{ and }
\cP_{1,F}\cap\{\fp\in\Pi_F\backslash\cP_F^{(d)}:\Frob_\fp=k\}
\text{ differ by a finite set.}
\end{equation}
On the other hand, since $H$ is abelian, the Chebotarev density theorem gives
$$
\delta_F(\{\fp\in\Pi_F\backslash\cP_F^{(d)}:\Frob_\fp=k\})=\frac{1}{|H|}.
$$
Since we have a partition $\Pi_F=\cP_{1,F}\sqcup \cP_{2^+,F}$,
and $\delta_F(\cP_{1,F})=1$ and $\delta_F(\cP_{2^+,F})=0$
by Lemma~\ref{L:degree1}, we have
\begin{equation}\label{E:delta2}
\delta_F(\cP_{1,F}\cap\{\fp\in\Pi_F\backslash\cP_F^{(d)}:\Frob_\fp=k\})=\frac{1}{|H|}.
\end{equation}
Combining \eqref{E:finite}, \eqref{E:delta1}, and~\eqref{E:delta2}, we conclude
\[
\delta_F(\Ord(C/F(t))\cap \cP_{1,F,k})=\frac{1}{|H|}.
\]
Since $\bigsqcup_{k\in \cK_\va}\cP_{1,F,k}$ is a disjoint union,
we have
\[
\delta_F\left(\Ord(C/F(t))\cap \bigsqcup_{k\in\cK_\va}\cP_{1,F,k}\right)
=\sum_{k\in\cK_\va}\delta_F\left(\Ord(C/F(t))\cap\cP_{1,F,k}\right)
=\frac{|\cK_\va|}{|H|}.
\]
But $|\cK_\va|\ge1$, so we obtain
\[
\underline{\delta}_F(\Ord(C/F(t)))\ge
\delta_F\left(\Ord(C/F(t))\cap \bigsqcup_{k\in\cK_\va}\cP_{1,F,k}\right)
=\frac{|\cK_\va|}{|H|}\ge \frac{1}{|H|}=\frac{1}{[L:F]}.
\]
This proves the first statement.

(2)
Let $\cP_{1,L,\ge 2d-1}$ denote the
subset of $\cP_{1,L}$ whose residue characteristic $p$ satisfies
$p\ge 2d-1$.

Let $\fP\in\cP_{1,L,\ge 2d-1}$.
Then $\kappa(\fP)=\F_p$, and $p\nmid d$.
Since $p\nmid d$, the reduction homomorphism
$\cO_{L,\fP}^\times \longrightarrow \kappa(\fP)^\times$
is injective on its prime-to-$p$ torsion subgroup.
Indeed, suppose that $u\equiv 1\pmod \fP$ and $u^n=1$, where $p\nmid n$.
Then $$0=u^n-1=(u-1)(1+u+\cdots+u^{n-1}),$$
while
\[
1+u+\cdots+u^{n-1}\equiv n \not\equiv 0 \pmod \fP,
\]
so the second factor is a unit in $\cO_{L,\fP}$, and hence $u=1$.
Since $\zeta_d\in\cO_L^\times$
has order $d$ and $p\nmid d$, its reduction also has order $d$ in
$\kappa(\fP)^\times = \F_p^\times$.
Hence
$d\mid p-1$, so
$p\equiv 1\pmod d$.
Theorem~\ref{T:main3} therefore shows that
$\fP\in\Ord(C_L/L(t))$.
This proves
\begin{equation}\label{E:include}
\cP_{1,L,\ge 2d-1}\subseteq \Ord(C_L/L(t)).
\end{equation}

By \eqref{E:include}, $\Ord(C_L/L(t))\cap\cP_{1,L}$ contains
$\cP_{1,L,\ge 2d-1}$.
Since $\cP_{1,L}\backslash \cP_{1,L,\ge 2d-1}$ is finite and $\delta_L(\cP_{1,L})=1$
by Lemma~\ref{L:degree1}, we have
\begin{equation}\label{E:fff}
\delta_L(\Ord(C_L/L(t))\cap \cP_{1,L})=\delta_L(\cP_{1,L})=1.
\end{equation}
Thus $\delta_L(\Ord(C_L/L(t)))=1$.
\end{proof}

\end{document}
